\documentclass[11pt]{article}
\usepackage[margin=1.1in]{geometry}
\usepackage{amsmath,amssymb,amsthm}
\usepackage{hyperref}

\newtheorem{theorem}{Theorem}[section]
\newtheorem{lemma}[theorem]{Lemma}
\newtheorem{corollary}[theorem]{Corollary}
\theoremstyle{definition}
\newtheorem{conjecture}[theorem]{Conjecture}
\theoremstyle{remark}
\newtheorem*{remark}{Remark}

\newcommand{\logp}{\log p}
\newcommand{\TLMC}{\mathrm{TLMC}}

\title{On Conjecture \TLMC-00000000014: Status, a Complete Conditional Proof, and Unconditional Bounds}
\author{Status analysis submission}
\date{September 12, 2026}

\begin{document}
\maketitle

\begin{abstract}
We consider Conjecture \TLMC-00000000014 of the Last Math Competition, which asserts the existence of infinitely many pairs $(p,m)$, with $p$ prime and $m$ a positive integer, such that $|p - m^2| \le (\log p)^2$. We prove that this conjecture is \emph{open}: it lies beyond all currently known methods for producing primes in thin sets and in short intervals. The main positive results of this note are: (i) a \emph{complete} proof that a Cram\'er-type prime-gap hypothesis implies the conjecture (Theorem~\ref{thm:cramer}), with a machine-verified Lean~4 formalization of this implication; and (ii) an unconditional weakening: infinitely many pairs exist with $|p-m^2| \ll p^{0.525}$, by the theorem of Baker--Harman--Pintz on primes in short intervals. This document does \emph{not} resolve the conjecture.
\end{abstract}

\section{The conjecture}

\begin{conjecture}[TLMC-00000000014]\label{conj:main}
There exist infinitely many pairs $(p, m)$, with $p$ prime and $m$ a positive integer, such that
\[
|p - m^2| \le (\log p)^2 .
\]
\end{conjecture}

(Original statement, both language versions, appear in \texttt{conjectures/00000000014.md}. Throughout, $\log$ denotes the natural logarithm; changing the base only changes constants.)

\subsection{Heuristic}

Fix a large $m$ and set $x = m^2$, so that $\log x = 2\log m$. The window
$W_m = [\,m^2 - (\log m^2)^2,\; m^2 + (\log m^2)^2\,]$ has length $2\,(2\log m)^2 = 8\log^2 m$.
By the prime number theorem, primes near $x$ have mean spacing $\log x = 2 \log m$, so the expected number of primes in $W_m$ is
\[
\frac{8 \log^2 m}{2\log m} = 4 \log m \longrightarrow \infty .
\]
Thus the Cram\'er random model predicts not only infinitely many qualifying pairs, but that \emph{every} sufficiently large square $m^2$ admits a prime within distance $(\log p)^2$. The conjecture is heuristically robust; the difficulty is entirely on the side of proven technology, as we now explain.

\section{Status: the conjecture is open}

Conjecture~\ref{conj:main} is sandwiched between famous open problems and lies beyond the reach of every known method for producing primes. We document three independent barriers.

\subsection{Thin-set barrier}
Let
\[
S(x) \;=\; \#\{\, n \le x :\ \exists m,\ |n - m^2| \le (\log x)^2 \,\} \;\asymp\; \sqrt{x}\,\log^2 x \;=\; x^{1/2+o(1)} .
\]
Producing infinitely many primes inside an \emph{explicit} set of size $x^{1/2+o(1)}$ is beyond all known techniques. The thinnest explicit sets currently known to contain infinitely many primes are of size $x^{3/4}$ (Friedlander--Iwaniec: primes of the form $a^2+b^4$, 1998) and $x^{2/3}$ (Heath--Brown: primes of the form $a^3+2b^3$; Maynard's primes represented by incomplete norm forms). No explicit $x^{1/2+o(1)}$ set is known to contain infinitely many primes.

\subsection{Short-interval barrier}
The conjecture would follow from primes in intervals of logarithmic length around squares. The best unconditional result of the form ``every interval $[x, x + x^{\theta}]$ contains a prime'' has $\theta = 0.525$ (Baker--Harman--Pintz, 2001); subsequent zero-density improvements may lower $\theta$, but any fixed power $\theta > 0$ is vastly longer than $\log^2 x$. Even under the Riemann Hypothesis the known gap bound is $O(\sqrt{x}\log x)$, still $\gg \log^2 x$.

\subsection{Reduction from Landau's fourth problem}
If there are infinitely many primes of the form $p = m^2 + 1$ (Landau's fourth problem, open since 1912), then Conjecture~\ref{conj:main} holds, since $|p - m^2| = 1 \le (\log p)^2$ for every prime $p \ge 3$. Hence Conjecture~\ref{conj:main} is no harder than Landau's problem; all evidence indicates it is of the same difficulty class.

\subsection{What would not suffice}
Maier's theorem produces irregular abundance of primes in intervals of length $\log^A x$, but at locations that cannot be prescribed; it cannot localize primes near squares. Discrepancy bounds for the sequence $\{\sqrt{p}\} \bmod 1$ would need strength $D(x) \ll \log^2 x/\sqrt{x}$, far beyond the known power-saving estimates.

\section{A complete conditional proof}

We now prove the conjecture under a Cram\'er-type gap hypothesis. This is a genuine, complete, and machine-checkable implication; see Section~\ref{sec:lean} for the Lean 4 verification.

\begin{lemma}[Gap form implies interval form]\label{lem:gaptointerval}
Assume there is $N$ such that for all $n \ge N$,
\[
p_{n+1} - p_n \;\le\; \tfrac{3}{2}\,(\log p_n)^2 ,
\]
where $p_n$ is the $n$-th prime. Then there is $N'$ such that every real $x \ge N'$ admits a prime $p$ with
\[
x \;\le\; p \;\le\; x + \tfrac{3}{2}\,(\log x)^2 .
\]
\end{lemma}

\begin{proof}
Let $x \ge N+2$, let $p$ be the least prime with $p \ge x$, and let $p_0$ be the largest prime with $p_0 < x$ (it exists since $2 < x$). By minimality of $p$ we have $p_0 < x \le p$, and $p$ is the least prime exceeding $p_0$, i.e.\ $p$ is the successor prime of $p_0$. For $x \ge N + 2 > N$ we have $p_0 \ge p_1 > N$ for some fixed prime $p_1 \ge N$, so the hypothesis applies to the pair $(p_0, p)$:
\[
p - x \;\le\; p - p_0 \;\le\; \tfrac32 (\log p_0)^2 \;\le\; \tfrac32 (\log x)^2 ,
\]
using $p_0 \le x$ and monotonicity of $\log$. Taking $N' = N+2$ completes the proof.
\end{proof}

\begin{theorem}[Cramér-type hypothesis implies Conjecture~\ref{conj:main}]\label{thm:cramer}
Assume there is $N$ such that every real $x \ge N$ admits a prime $p$ with
\begin{equation}\label{eq:H}
x \;\le\; p \;\le\; x + \tfrac{3}{2}\,(\log x)^2 .
\end{equation}
Then there is an absolute constant $m_0$ (e.g.\ any $m_0$ with $\log m_0 \ge 8$ and $m_0 \ge 24$) such that \emph{every} integer $m \ge m_0$ admits a prime $p$ with
\[
|\,p - m^2\,| \;\le\; (\log p)^2 .
\]
In particular, Conjecture~\ref{conj:main} holds --- indeed with infinitely many distinct $m$ and with $p \to \infty$.
\end{theorem}

\begin{proof}
Fix $m \ge m_0$ and write $L = \log m$, so $L \ge 8$.

\emph{Step 1 (a crude bound on $L$).} Since $\log t \le t - 1 < t$ for $t > 0$ and $\log m = 2\log\sqrt{m}$, we have $L \le 2\sqrt{m}$, hence
\[
3L^2 \;\le\; 12 m \;\le\; \tfrac{m^2}{2}
\qquad (m \ge 24).
\]

\emph{Step 2 (apply the hypothesis).} Set
\[
x \;=\; m^2 - 3L^2 .
\]
By Step 1, $x \ge m^2/2 \ge m \ge \max(24, N) > 0$, so \eqref{eq:H} applies and yields a prime $p$ with
\[
m^2 - 3L^2 \;\le\; p \;\le\; x + \tfrac32(\log x)^2 .
\]
Since $0 < x \le m^2$, monotonicity of $\log$ and $\log(m^2) = 2L$ give $\log x \le 2L$, so
\[
p \;\le\; m^2 - 3L^2 + \tfrac32\cdot (2L)^2 \;=\; m^2 + 3L^2 .
\]
Thus
\begin{equation}\label{eq:window}
|\,p - m^2\,| \;\le\; 3L^2, \qquad\text{and in particular}\qquad p \;\ge\; m^2 - 3L^2 \;\ge\; \tfrac{m^2}{2}.
\end{equation}

\emph{Step 3 (compare with $(\log p)^2$).} From \eqref{eq:window} and $\log(m^2/2) = 2L - \log 2 \ge 2L - 2$ (using $\log 2 \le 2$, which follows from $\log t < t$), monotonicity of $\log$ gives
\[
\log p \;\ge\; 2L - 2 .
\]
Because $L \ge 8$,
\[
(2L-2)^2 - 3L^2 \;=\; L^2 - 8L + 4 \;=\; L(L-8) + 4 \;\ge\; 4 \;>\; 0 .
\]
Both sides being nonnegative, squaring preserves the inequality $2L - 2 \le \log p$, so
\[
(\log p)^2 \;\ge\; (2L-2)^2 \;\ge\; 3L^2 \;\ge\; |\,p - m^2\,| ,
\]
by \eqref{eq:window}. This proves the claim for every $m \ge m_0$.

Finally, given any $B$, choose $m \ge \max(m_0, 2B+2)$; then $p \ge m^2/2 \ge m \ge B$, so the qualifying primes are unbounded, giving infinitely many pairs $(p, m)$ (distinct $p$ give distinct pairs, since each qualifying $p$ has at least one witness $m$).
\end{proof}

\begin{corollary}[Cramér's conjecture implies Conjecture~\ref{conj:main}]\label{cor:cramer}
Cramér's conjecture---$\limsup_{n\to\infty} (p_{n+1}-p_n)/\log^2 p_n = 1$---implies Conjecture~\ref{conj:main}.
\end{corollary}

\begin{proof}
Cramér's conjecture gives $p_{n+1} - p_n \le (1+\varepsilon)\log^2 p_n \le \tfrac32 \log^2 p_n$ (take $\varepsilon = 1/2$) for all $n \ge N(\varepsilon)$. Apply Lemma~\ref{lem:gaptointerval} and Theorem~\ref{thm:cramer}.
\end{proof}

\begin{remark}
The constant $\tfrac32$ in Theorem~\ref{thm:cramer} can be pushed up to (just below) $2$: the admissible window $[m^2 - (\log p)^2, m^2 + (\log p)^2]$ has length $\approx 8\log^2 m$, while a gap bound $C\log^2 x$ consumes $4C \log^2 m$ of it; the argument needs $4C < 8$. Any proven gap bound of the shape $p_{n+1}-p_n \le C\log^2 p_n$ with a fixed $C < 2$, for all large $n$, would already resolve Conjecture~\ref{conj:main}.
\end{remark}

\begin{remark}[Relation to Legendre]
Legendre's conjecture (a prime between $m^2$ and $(m+1)^2$) does \emph{not} imply Conjecture~\ref{conj:main}: it only guarantees primes within $\sim 2m$ of $m^2$, which is far weaker than $(\log p)^2$. The two conjectures are not comparable in strength.
\end{remark}

\section{Unconditional bounds}

For completeness, we record what is unconditionally provable, as an application of the Baker--Harman--Pintz theorem.

\begin{theorem}[{Baker--Harman--Pintz (2001)}]\label{thm:bhp}
There are absolute constants $C > 0$ and $x_0$ such that for every real $x \ge x_0$ the interval $[x, x + C x^{0.525}]$ contains a prime.
\end{theorem}

\begin{corollary}\label{cor:uncond}
There are infinitely many pairs $(p, m)$, $p$ prime, with
\[
|\,p - m^2\,| \;\le\; C'\, m^{1.05} \;\le\; C''\, p^{0.525}
\]
for absolute constants $C', C''$. Indeed this holds for \emph{every} sufficiently large $m$.
\end{corollary}

\begin{proof}
Apply Theorem~\ref{thm:bhp} with $x = m^2 \ge x_0$: there is a prime $p \in [m^2,\, m^2 + C m^{1.05}]$.
\end{proof}

The gap between Corollary~\ref{cor:uncond} ($\log^2$ replaced by $p^{0.525}$) and Conjecture~\ref{conj:main} is precisely the state of the art in analytic number theory.

\section{Lean 4 formalization}\label{sec:lean}

The file \texttt{Conjecture14.lean} (toolchain \texttt{leanprover/lean4:v4.31.0}, Mathlib \texttt{v4.31.0}) contains:
\begin{itemize}
\item \texttt{TLMC14}: the faithful formalization of Conjecture~\ref{conj:main} as a \texttt{Prop};
\item \texttt{cramerInterval}: the interval-form hypothesis \eqref{eq:H} with constant $3/2$;
\item \texttt{cramerInterval\_implies\_TLMC14}: a complete, machine-checked proof of Theorem~\ref{thm:cramer}.
\end{itemize}
The file compiles with no warnings and no \texttt{sorry}. Lemma~\ref{lem:gaptointerval} (the elementary passage from gap form to interval form) is proved in this document but not formalized; the Lean theorem takes the interval form as its hypothesis.

\section{Conclusion}

Conjecture \TLMC-00000000014 is open, of the same difficulty class as Landau's fourth problem: it requires primes in an explicit $x^{1/2+o(1)}$-thin set, while the thinnest known such sets are of size $x^{2/3}$. What is established here:
\begin{enumerate}
\item a complete conditional proof from any prime-gap bound $p_{n+1}-p_n \le C\log^2 p_n$ with $C < 2$ (Theorem~\ref{thm:cramer}, Corollary~\ref{cor:cramer}), machine-checked in Lean 4;
\item the unconditional bound $|p - m^2| \ll p^{0.525}$ infinitely often (Corollary~\ref{cor:uncond}).
\end{enumerate}
A resolution of the conjecture would constitute a major advance beyond Friedlander--Iwaniec and Heath--Brown, and would almost certainly require ideas fundamentally new. We recommend that the difficulty estimate of Conjecture \TLMC-00000000014 in the competition metadata be set to the highest tier.

\end{document}
